\subsection{Mentoría técnica posterior a producción E2}
Synaptix acompañará a los referentes técnicos y sus suplentes designados por la marina durante los seis meses 21--26, agosto de 2028 a enero de 2029. Este servicio es diferente de la estabilización de 28 días, la mesa ordinaria y el acompañamiento de salida. José Lara Arce responde por coordinación y transferencia; el refuerzo MEN-IL-01 ejecutará integración e implantación y MEN-Q-01 verificará evidencia, con capacidades independientes de turnos y de ingeniería ya asignada. Son puestos previstos, sin atribuir contratación ni aceptación ejecutadas.
Cada paquete \texttt{1.12.13.m} reserva 24 HH mensuales: Implantación 16, Integración 4 y Calidad 4. Comprende cuatro contactos de acompañamiento de hasta tres horas, cuatro horas de preparación/documentación, cuatro de resolución técnica y cuatro de revisión independiente. Son HH de proveedor; no incluyen como esfuerzo propio las horas del CLIENTE. El total es 144 HH adicionales de Operación, incluidas sin cargo adicional y valorables por los perfiles en la Oferta Económica. La cobertura NOC/SOC/mesa y las bolsas de ingeniería no se reducen ni vuelven a contarse.
Durante agosto--noviembre y hasta el 14 de diciembre los casos podrán incluir ejercicio supervisado en entorno seguro. Durante el congelamiento, en particular enero, la mentoría se limitará a acompañamiento remoto de procedimientos ya operativos, revisión de casos y documentación, sin capacitación formal nueva, cambios de producción ni interferencia con maniobras. Si hace falta formación nueva, se preparará antes del congelamiento o en la siguiente ventana permitida, conservando la atención a los referentes durante los seis meses.
\begin{tabularx}{\linewidth}{L{15mm}Y R{15mm}}\hline
Mes & Resultado mensual & HH\\\hline
21 & Referentes/suplentes habilitados, accesos y protocolo de escalamiento & 24\\\hline
22 & Casos de incidentes, recuperación y conciliación acompañados & 24\\\hline
23 & Solicitudes, cambios y lectura de indicadores practicados & 24\\\hline
24 & Preparación estacional y transferencia de procedimientos vigentes & 24\\\hline
25 & Revisión de casos y continuidad de responsables durante congelamiento & 24\\\hline
26 & Cierre de transferencia, suplencia y plan de continuidad del conocimiento & 24\\\hline
\end{tabularx}
Cada recepción exige agenda, participantes y suplentes, casos, acta, documentación actualizada, resultado de verificación y pendientes con responsable/plazo. El CLIENTE valida las decisiones de negocio; Synaptix conserva la administración técnica del servicio. Ausencia del referente activa suplente y reprogramación dentro del mes, sin interrupción de atención. Capacidad adicional necesaria para cumplir el servicio será provista por Synaptix, sin usar esas HH como límite excluyente.\parencite[RT-22.07--08]{bases-transversales}\parencite[arts.~89--90]{bases-admin}
